
\documentclass[11pt]{article}
\usepackage[a4paper,margin=1in]{geometry}
\usepackage[T1]{fontenc}
\usepackage[utf8]{inputenc}
\usepackage{lmodern}
\usepackage{microtype}
\usepackage{amsmath,amssymb,mathtools}
\usepackage{amsthm}
\usepackage{enumitem}
\usepackage{tikz-cd}
\usepackage[hidelinks]{hyperref}
\setlength{\parindent}{0pt}
\setlength{\parskip}{0.55em}
\newcommand{\Aone}{\mathbb A^1}
\newcommand{\A}{\mathbb A}
\newcommand{\Z}{\mathbb Z}
\newcommand{\Q}{\mathbb Q}
\newcommand{\Gm}{\mathbb G_m}
\newcommand{\Et}{\mathrm{Et}}
\newcommand{\et}{\mathrm{et}}
\newcommand{\Nis}{\mathrm{Nis}}
\newcommand{\Zar}{\mathrm{Zar}}
\newcommand{\Sm}{\mathrm{Sm}}
\newcommand{\SmCor}{\mathrm{SmCor}}
\newcommand{\Spec}{\operatorname{Spec}}
\newcommand{\Hom}{\operatorname{Hom}}
\newcommand{\Ker}{\operatorname{Ker}}
\newcommand{\Coker}{\operatorname{Coker}}
\newcommand{\Tot}{\operatorname{Tot}}
\newcommand{\forget}{\operatorname{forget}}
\newcommand{\Sing}{C_*}
\newcommand{\D}{D}
\newcommand{\Ho}{\mathcal H}
\newcommand{\Hob}{\mathcal H_\bullet}
\newcommand{\HoA}{\mathcal H_{\bullet,A^1}}
\newcommand{\Ch}{\mathrm{Ch}}
\newcommand{\sSh}{\mathrm{sSh}}
\newcommand{\Sh}{\mathrm{Sh}}
\newcommand{\pt}{*}
\newcommand{\ZZ}{\widetilde{\Z}}
\newcommand{\HH}{\mathrm H}
\newcommand{\red}{\mathrm{red}}
\newtheoremstyle{deligne}{0.5em}{0.5em}{\itshape}{}{}{.}{0.5em}{\bfseries #1\thmnote{ #3}}
\theoremstyle{deligne}
\newtheorem*{theorem*}{Théorème}
\newtheorem*{proposition*}{Proposition}
\newtheorem*{lemma*}{Lemme}
\newtheorem*{corollary*}{Corollaire}
\newtheorem*{definition*}{Définition}
\newtheorem*{remark*}{Remarque}
\newtheorem*{example*}{Exemple}
\newcommand{\heading}[1]{\vspace{0.7em}\noindent{\bfseries #1}\quad}
\newenvironment{proofdel}{\noindent\textit{Preuve.}}{\hfill$\square$}
\newenvironment{proofnamed}[1]{\noindent\textit{#1.}}{}

\newcommand{\uHom}{\underline{\Hom}}
\newcommand{\sk}{\operatorname{sk}}
\newcommand{\Wr}{\operatorname{Wr}}
\newcommand{\colim}{\operatorname{colim}}
\newcommand{\diag}{\Delta}
\newcommand{\inj}{\mathrm{inj}}
\newcommand{\op}{\mathrm{op}}
\newcommand{\limi}{\operatorname{lim}}
\newcommand{\Ex}{\operatorname{Ex}}
\newcommand{\Id}{\operatorname{Id}}
\newcommand{\cyl}{\operatorname{cyl}}

\begin{document}

\begin{center}
{\Large\bfseries Leçons de Voevodsky sur la cohomologie motivique\\
2000/2001}\\[0.5em]
{\large Pierre Deligne}
\end{center}

\section*{Suite de 4.1}

\begin{proofdel}
Le carré précédent est cocartésien, avec \(f\) ouvert. En particulier \(f\) est un monomorphisme; il s'ensuit que \(f'\) est un monomorphisme et que le carré est aussi cartésien. Le morphisme \(F'\amalg G\to G'\) est surjectif. Le tiré en arrière de \(f'\) par \(f':F'\to G'\) est un isomorphisme, puisque \(f'\) est un monomorphisme, et il est donc ouvert. Le tiré en arrière de \(f'\) par \(G\to G'\) n'est autre que \(f\), ouvert par hypothèse. Il en résulte que \(f'\) est ouvert.
\end{proofdel}

Nous fixons maintenant une classe \(\mathcal E\) de plongements ouverts \(U\to X\) dans \((QP/G)\). Nous exigeons les propriétés de stabilité suivantes.

\begin{definition*}[Conditions 34]
\begin{enumerate}[label=\arabic*.]
\item Si \(U\to U'\to X\) sont des plongements ouverts et si \(U\to X\) appartient à \(\mathcal E\), alors \(U'\to X\) appartient à \(\mathcal E\).
\item Si \(U\to X\) est un plongement ouvert de \(\mathcal E\), et si \(f:X'\to X\) est étale, induisant un isomorphisme de \(f^{-1}(Z)\) sur \(Z\), où \(Z\) est le complémentaire de \(U\) dans \(X\), alors \(f^{-1}(U)\to X'\) appartient à \(\mathcal E\).
\end{enumerate}
\end{definition*}

Les classes \(\mathcal E\) que nous aurons à considérer sont les suivantes: les plongements ouverts \(U\to X\) avec \(X\) lisse; les plongements ouverts \(U\to X\) avec \(X\) lisse et l'action de \(G\) libre sur \(X-U\), c'est-à-dire \(X\) réunion de \(U\) et de l'ouvert où l'action de \(G\) est libre; et, lorsqu'on travaille dans \((Sch/S)\), tous les plongements ouverts.

\begin{definition*}[7]
Un morphisme \(f:F\to G\) est \(\mathcal E\)-solide s'il est un composé
\[
F=F_0\longrightarrow F_1\longrightarrow \cdots\longrightarrow F_n=G
\]
où chaque \(F_i\to F_{i+1}\) est déduit par somme amalgamée d'un \(h_U\to h_X\), avec \(U\subset X\) dans \(\mathcal E\). Un faisceau \(F\) est solide si \(\emptyset\to F\) est \(\mathcal E\)-solide. Dans le contexte pointé, un faisceau pointé est \(\mathcal E\)-solide si le morphisme \(\pt\to F\) est \(\mathcal E\)-solide.
\end{definition*}

\begin{example*}[1]
Pour \(U\) ouvert dans \(X\), soit \(h_X/h_U\) le faisceau \(h_X\) où le sous-faisceau \(h_U\) est contracté au point-base. Si \(U\to X\) est dans \(\mathcal E\), alors \(\pt\to h_X/h_U\) est solide: c'est le poussé en avant de \(h_U\to h_X\) par \(h_U\to\pt\). Les espaces de Thom sont de cette forme: à partir d'un fibré vectoriel \(E\) sur \(T\), on contracte, dans l'espace total de \(E\), le complément de la section nulle en un point.
\end{example*}

La classe des morphismes solides est la plus petite classe, stable par composition et sommes amalgamées, qui contient tous les \(h_U\to h_X\), \(U\subset X\) dans \(\mathcal E\). Par la proposition 33, les morphismes solides sont ouverts.

Pour un monomorphisme \(F\to G\) de faisceaux, on définit \(G/F\) comme le faisceau pointé obtenu en contractant \(F\) en un point. On a donc un carré cocartésien
\[
\begin{tikzcd}[column sep=2.6em,row sep=1.6em]
F \arrow[r] \arrow[d] & G \arrow[d]\\
\pt \arrow[r] & G/F .
\end{tikzcd}
\]
Par transitivité des sommes amalgamées, tout diagramme cocartésien
\[
\begin{tikzcd}[column sep=2.6em,row sep=1.6em]
F \arrow[r] \arrow[d] & G \arrow[d]\\
F' \arrow[r] & G'
\end{tikzcd}
\]
induit un isomorphisme \(G/F\simeq G'/F'\).

\begin{proposition*}[35]
Un morphisme de faisceaux \(f:F\to G\) est \(\mathcal E\)-solide si et seulement s'il est un composé
\[
F=F_0\longrightarrow F_1\longrightarrow \cdots\longrightarrow F_n=G
\]
de monomorphismes tel que chaque \(F_{i+1}/F_i\) soit isomorphe à un \(h_X/h_U\), avec \(U\subset X\) dans \(\mathcal E\).
\end{proposition*}

\begin{proofdel}
Si un morphisme \(F\to G\) est déduit par somme amalgamée de \(U\to X\), alors \(G/F\) est isomorphe à \(h_X/h_U\). Ceci donne le sens «seulement si».

Réciproquement, supposons d'abord que l'on ait un carré cocartésien
\[
\begin{tikzcd}[column sep=2.8em,row sep=1.8em]
F \arrow[r] \arrow[d] & G \arrow[d]\\
\pt \arrow[r] & h_X/h_U
\end{tikzcd}
\tag{35.1}
\]
et que le morphisme naturel \(h_X\to h_X/h_U\) se relève à \(G\). Alors \(F\to G\) est obtenu par somme amalgamée à partir de \(h_U\to h_X\). En effet, (35.1) est cartésien aussi bien que cocartésien; on a donc un carré cartésien
\[
\begin{tikzcd}[column sep=2.8em,row sep=1.8em]
h_U \arrow[r] \arrow[d] & h_X \arrow[d]\\
F \arrow[r] & G .
\end{tikzcd}
\]
Si \(G'\) est déduit de \(F\) par somme amalgamée de \(h_U\to h_X\),
\[
\begin{tikzcd}[column sep=2.3em,row sep=1.6em]
h_U \arrow[r] \arrow[d] & h_X \arrow[d] \arrow[r,equal] & h_X \arrow[d]\\
F \arrow[r] & G' \arrow[r] & G ,
\end{tikzcd}
\]
alors \(G'/F\simeq G/F\), et il s'ensuit que \(G'\simeq G\).

Il suffit plus généralement d'avoir un morphisme étale \(v:\widetilde X\to X\), induisant un isomorphisme de \(\widetilde X-v^{-1}(U)\) sur \(X-U\), et un relèvement de \(h_{\widetilde X}\to h_X/h_U\) à \(G\). Si \(\widetilde U=v^{-1}(U)\), alors \(h_{\widetilde X}/h_{\widetilde U}\to h_X/h_U\) est un isomorphisme. Cela se vérifie sur les foncteurs fibres définis par les schémas henséliens \(G\)-locaux: un morphisme \(T\to X\) qui n'arrive pas dans \(U\) se relève de façon unique à \(\widetilde X\). Sous ces hypothèses, \(F\to G\) est une somme amalgamée de \(h_{\widetilde U}\to h_{\widetilde X}\). Par la seconde propriété de stabilité de \(\mathcal E\), le plongement \(\widetilde U\to\widetilde X\) appartient à \(\mathcal E\).

On réduit le cas général à celui-ci. Soit \(f:F\to G\) un monomorphisme tel que \(G/F\simeq h_X/h_U\), \(U\to X\) dans \(\mathcal E\). Le carré cocartésien
\[
\begin{tikzcd}[column sep=2.8em,row sep=1.8em]
F \arrow[r] \arrow[d] & G \arrow[d]\\
\pt \arrow[r] & G/F
\end{tikzcd}
\tag{35.2}
\]
donne un épimorphisme \(\pt\amalg G\to h_X/h_U\). La section naturelle de \(h_X/h_U\) sur \(X\) peut donc être relevée localement soit à \(\pt\), soit à \(G\). Ainsi, pour une filtration
\[
\emptyset=C_{n+1}\subset C_n\subset\cdots\subset C_1\subset C_0=X
\]
par sous-schémas fermés équivariants, il existe des morphismes étales \(T_i\to X\), munis de sections équivariantes au-dessus de \(C_i-C_{i+1}\), et des relèvements de \(h_{T_i}\to h_X/h_U\) à \(\pt\) ou à \(G\). On peut commencer par \(X_0=U\), où le relèvement est à \(\pt\), puis rétrécir les \(T_i\) de sorte que les relèvements suivants soient à \(G\) et induisent des isomorphismes au-dessus des strates.

Comme \(F\to G\) est un monomorphisme, (35.2) est aussi cartésien. En tirant en arrière le composé
\[
\pt\longrightarrow h_{X_1}/h_U\longrightarrow h_{X_2}/h_U\longrightarrow\cdots\longrightarrow h_X/h_U
\]
on obtient une factorisation de \(F\to G\) dont chaque quotient est de la forme \(h_{T_i}/h_{T_i'}\), et où les morphismes \(h_{T_i}\to h_X/h_U\) se relèvent aux faisceaux intermédiaires. Chaque étape est donc une somme amalgamée d'un plongement ouvert de \(\mathcal E\), et la solidité s'ensuit.
\end{proofdel}

\begin{remark*}[1]
Autrement dit, \(F\to G\) est \(\mathcal E\)-solide si et seulement si le faisceau pointé \(G/F\) est une extension itérée de faisceaux \(h_X/h_U\), avec \(U\to X\) dans \(\mathcal E\): il existe des morphismes
\[
\pt=E_0\to E_1\to\cdots\to E_n=G/F
\]
avec chaque quotient \(E_{i+1}/E_i\) de la forme \(h_X/h_U\).
\end{remark*}

\begin{proposition*}[36]
Si \(f:F\to G\) est ouvert et si \(G\) est \(\mathcal E\)-solide, alors \(f\) est \(\mathcal E\)-solide.
\end{proposition*}

\begin{proofdel}
Nous disons «solide» pour \(\mathcal E\)-solide. Soit \((*)\) la propriété d'un faisceau \(G\): tout morphisme ouvert \(F\to G\) est solide. Si \(G\) est solide, il apparaît au bout d'une chaîne
\[
\emptyset=G_0\to G_1\to\cdots\to G_n=G
\]
dans laquelle chaque \(G_i\to G_{i+1}\) est une somme amalgamée d'un \(h_U\to h_X\), \(U\subset X\) dans \(\mathcal E\). On prouve par récurrence que \(G_i\) vérifie \((*)\).

Pour \(n=0\), \(G_1=h_X\) est représentable et \(\emptyset\to X\) appartient à \(\mathcal E\). Si \(F\to h_X\) est ouvert, il est de la forme \(h_U\to h_X\), pour \(U\) ouvert dans \(X\), donc solide par la condition 34(1).

Il reste l'étape de récurrence. Dans un carré cocartésien
\[
\begin{tikzcd}[column sep=2.8em,row sep=1.8em]
h_U \arrow[r] \arrow[d] & h_X \arrow[d]\\
G' \arrow[r] & G
\end{tikzcd}
\tag{36.1}
\]
le morphisme \(h_U\to h_X\) est un monomorphisme; le carré est donc cartésien aussi bien que cocartésien. Soit \(F\to G\) ouvert. En tirant (36.1) en arrière par \(F\to G\), on obtient un carré cartésien et cocartésien
\[
\begin{tikzcd}[column sep=2.8em,row sep=1.8em]
h_V \arrow[r] \arrow[d] & h_Y \arrow[d]\\
F' \arrow[r] & F
\end{tikzcd}
\tag{36.2}
\]
où \(Y\) est ouvert dans \(X\) et \(V=U\cap Y\). Le diagramme
\[
\begin{tikzcd}[column sep=2.6em,row sep=1.55em]
F' \arrow[r] \arrow[d] & F \arrow[r] \arrow[d] & h_Y/h_V \arrow[d]\\
G' \arrow[r] \arrow[d] & G \arrow[r] \arrow[d] & h_X/h_U \arrow[d]\\
G'/F' \arrow[r] & G/F \arrow[r] & h_X/(h_U\cup h_Y)
\end{tikzcd}
\]
exprime \(G/F\) comme une extension de \(h_X/h_{U\cup Y}\) par \(G'/F'\). Par l'hypothèse de récurrence appliquée à \(G'\), et par la remarque 1, on en déduit que \(F\to G\) est solide. Ici \(U\cup Y\to X\) appartient à \(\mathcal E\).
\end{proofdel}

\begin{proposition*}[37]
Le tiré en arrière d'un morphisme solide par un morphisme ouvert est solide. En particulier, si \(F\to G\) est ouvert et si \(G\) est solide, alors \(F\) est solide.
\end{proposition*}

\begin{proofdel}
Comme le tiré en arrière d'un morphisme ouvert est ouvert, il suffit de vérifier l'assertion pour un morphisme solide obtenu comme somme amalgamée d'un \(h_U\subset h_X\):
\[
\begin{tikzcd}[column sep=2.8em,row sep=1.8em]
h_U \arrow[r] \arrow[d] & h_X \arrow[d]\\
G' \arrow[r] & G .
\end{tikzcd}
\]
En tirant par un morphisme ouvert \(F\to G\), on obtient un carré cocartésien
\[
\begin{tikzcd}[column sep=2.8em,row sep=1.8em]
h_{U'} \arrow[r] \arrow[d] & h_{X'} \arrow[d]\\
F' \arrow[r] & F
\end{tikzcd}
\]
avec \(U'\) ouvert dans \(U\) et \(X'\) ouvert dans \(X\). Ceci montre que \(F'\to F\) est solide.
\end{proofdel}

Supposons maintenant données deux classes \(\mathcal E\) et \(\mathcal E'\) de plongements ouverts vérifiant les conditions 34. On définit \(\mathcal E\times\mathcal E'\) comme la plus petite classe, stable par les conditions 34, contenant les plongements ouverts
\[
(U\times X')\cup(X\times U')\subset X\times X',
\]
pour \(U\subset X\) dans \(\mathcal E\) et \(U'\subset X'\) dans \(\mathcal E'\).

\begin{example*}[2]
Si \(\mathcal E\) est la classe de tous les plongements ouverts et si \(\mathcal E'\) est la classe des \(U'\subset X'\) tels que \(G\) agisse librement hors de \(U'\), alors \(\mathcal E\times\mathcal E'=\mathcal E'\).
\end{example*}

\begin{proposition*}[38]
Si les faisceaux pointés \(F\) et \(F'\) sont respectivement \(\mathcal E\)- et \(\mathcal E'\)-solides, alors \(F\wedge F'\) est \(\mathcal E\times\mathcal E'\)-solide.
\end{proposition*}

\begin{proofdel}
Par hypothèse, \(F\) est une extension itérée, au sens de la remarque 1, de faisceaux pointés \(h_{X_i}/h_{U_i}\), \(U_i\subset X_i\) dans \(\mathcal E\). De même \(F'\) est construit à partir de \(h_{X'_j}/h_{U'_j}\), \(U'_j\subset X'_j\) dans \(\mathcal E'\). Le smash-produit \(F\wedge F'\) est alors une extension itérée, par exemple dans l'ordre lexicographique, des
\[
(h_{X_i}/h_{U_i})\wedge(h_{X'_j}/h_{U'_j})
\simeq
h_{X_i\times X'_j}\big/
h_{(U_i\times X'_j)\cup(X_i\times U'_j)}.
\]
Il est donc \(\mathcal E\times\mathcal E'\)-solide.
\end{proofdel}

\begin{definition*}[8]
Un morphisme est dit ind-solide relativement à \(\mathcal E\) s'il est une limite inductive filtrante de morphismes \(\mathcal E\)-solides.
\end{definition*}

\heading{Exercice 39.}
On prend \(G\) trivial. Une section sur \(T\) d'une somme amalgamée
\[
\begin{tikzcd}[column sep=2.8em,row sep=1.8em]
h_U \arrow[r] \arrow[d] & h_X \arrow[d]\\
F \arrow[r] & G
\end{tikzcd}
\]
se décrit comme suit. Pour un ouvert \(V\) de \(T\) et une section \(s\) de \(F\) sur \(V\), considérons sur le petit site de Nisnevich \(T_{\Nis}\) le préfaisceau \(a(V,s)\) qui envoie \(Y\to T\) sur l'ensemble des morphismes \(f:Y\to X\) tels que \(f^{-1}(U)=Y^{-1}(V)\) et dont la restriction soit compatible avec \(s\). Une section de \(G\) sur \(T\) est donnée par:
\begin{enumerate}[label=\arabic*.]
\item un ouvert \(V\) de \(T\);
\item une section \(s\) de \(F\) sur \(V\);
\item une section de \(i_*(a(V,s)_{\Nis})\) sur \(T-V\), où \(i:T-V\to T\) est l'immersion fermée.
\end{enumerate}

\heading{Exercice 40.}
Avec les notations de l'exercice 39, si \(F\) est un faisceau pour la topologie étale, alors \(G\) l'est aussi. Pour \((V,s)\) fixé, la donnée (3) forme un faisceau étale. Cela se vérifie au moyen du critère usuel permettant de reconnaître qu'un faisceau de Nisnevich est étale: pour \(x\in T\) et pour une extension finie séparable \(K/k(x)\), après extension du corps résiduel à partir de l'hensélisé \(\mathcal O^h_{T,x}\), le foncteur correspondant des voisinages étales doit être un faisceau étale sur \(\Spec k(x)\).

\heading{Exercice 41.}
Il résulte des exercices 39 et 40 que si \(f:F\to G\) est ind-solide et si \(F\) est étale, alors \(G\) est étale. En particulier, les faisceaux solides, ainsi que les faisceaux pointés solides, sont des faisceaux étales.

\begin{remark*}[2]
Le même formalisme des morphismes ouverts et solides vaut dans le site de tous les schémas de type fini sur \(S\), muni de la topologie étale.
\end{remark*}

\subsection*{4.2 Un critère de préservation des équivalences locales}

Nous travaillons avec les faisceaux pointés sur \(QP/G\). Le but de cette section est de prouver le résultat suivant.

\begin{theorem*}[6]
Soit \(a\) un foncteur des faisceaux pointés vers les ensembles pointés. Supposons que \(a\) commute à toutes les colimites, et que, pour tout plongement ouvert \(U\to X\), le morphisme
\[
a((h_U)_+)\longrightarrow a((h_X)_+)
\]
soit un monomorphisme. Si \(f:F_\bullet\to G_\bullet\) est une équivalence locale et si les termes \(F_n\) et \(G_n\) sont ind-solides pointés, alors \(a(f)\) est une équivalence faible.
\end{theorem*}

Supposons que
\[
Q=\begin{tikzcd}[column sep=2.5em,row sep=1.6em,baseline=(current bounding box.center)]
B \arrow[r] \arrow[d] & Y \arrow[d]\\
A \arrow[r] & X
\end{tikzcd}
\]
soit un carré distingué supérieur. En ajoutant un point-base, on obtient \(Q_+\). Le morphisme \(K_Q\to X_+\) est une équivalence locale. Vérifions que \(a(K_Q)\to a(X_+)\) est une équivalence faible. Comme \(a\) commute aux coproduits, ce morphisme se déduit du carré commutatif
\[
\begin{tikzcd}[column sep=2.6em,row sep=1.8em]
a((h_B)_+) \arrow[r] \arrow[d] & a((h_Y)_+) \arrow[d]\\
a((h_A)_+) \arrow[r] & a((h_X)_+)
\end{tikzcd}
\]
par la même construction qu'en (23.3). Le carré est cocartésien parce que \(Q\) l'est. La flèche horizontale supérieure étant un monomorphisme, le carré est homotopiquement cocartésien, d'où l'assertion. Comme \(a\) commute aux colimites, ce cas particulier implique le lemme suivant.

\begin{lemma*}[15]
Si \(f\) appartient à la \(\bar\Delta\)-clôture des morphismes \((K_Q)_+\to X_+\) ci-dessus, alors \(a(f)\) est une équivalence faible.
\end{lemma*}

\begin{proposition*}[42]
Preuve du théorème 6.
\end{proposition*}

\begin{proofdel}
Pour tout faisceau pointé \(F\), le morphisme \(S(F)\to F\) est une équivalence locale: pour tout \(X\in QP/G\) connexe, \(S(F)(X)\to F(X)\) est une équivalence faible par la construction 29. Ainsi, pour une équivalence locale \(f:F\to G\), le carré
\[
\begin{tikzcd}[column sep=2.8em,row sep=1.8em]
S(F) \arrow[r] \arrow[d] & S(G) \arrow[d]\\
F \arrow[r] & G
\end{tikzcd}
\]
est un carré commutatif d'équivalences locales. Par les lemmes 15 et 12, \(a(S(f))\) est une équivalence faible. Il reste à montrer que \(a(S(F))\to a(F)\) est une équivalence faible, et de même pour \(G\). Comme \(a\) et \(S\) commutent aux colimites filtrantes, le cas ind-solide se ramène au cas solide; par la définition inductive de la solidité, il suffit de prouver le lemme suivant.
\end{proofdel}

\begin{lemma*}[16]
Soit \(U\to X\) un plongement ouvert. Supposons que dans un carré cartésien de faisceaux pointés
\[
\begin{tikzcd}[column sep=2.8em,row sep=1.8em]
(h_U)_+ \arrow[r] \arrow[d] & (h_X)_+ \arrow[d]\\
F \arrow[r] & G
\end{tikzcd}
\]
le faisceau \(F\) soit tel que \(a(S(F))\to a(F)\) soit une équivalence faible. Alors la même propriété vaut pour \(G\).
\end{lemma*}

\begin{proofdel}
Considérons le carré cocartésien
\[
\begin{tikzcd}[column sep=2.8em,row sep=1.8em]
S((h_U)_+) \arrow[r] \arrow[d] & S((h_X)_+) \arrow[d]\\
S(F) \arrow[r] & S(G).
\end{tikzcd}
\]
La flèche horizontale supérieure est un monomorphisme. Il s'ensuit que ce carré est point par point homotopiquement cocartésien, et que \(S\to S(G)\) est une équivalence locale. Le foncteur \(i^*i_*\) de la construction 29 transforme un monomorphisme en une coprojection de la forme \(F\to F\amalg\coprod(h_{U_\alpha})_+\). Par conséquent chaque \(S_n\) est de la forme \((\coprod h_{U_\alpha})_+\), et, par les lemmes 15 et 12, \(a(S)\to a(S(G))\) est une équivalence faible. Il reste à montrer que \(a(S)\to a(G)\) est une équivalence faible.

Appliquons \(a\) au morphisme du carré que l'on vient d'écrire vers le carré cartésien initial. Par la construction 29, les deux morphismes
\[
S((h_U)_+)\to(h_U)_+,\qquad S((h_X)_+)\to(h_X)_+
\]
sont des équivalences d'homotopie et le restent après application de \(a\). Nous avons supposé que \(a(S(F))\to a(F)\) était une équivalence faible. Comme \(a\) envoie les deux carrés sur des carrés cocartésiens dont les flèches horizontales supérieures sont des monomorphismes, il s'ensuit que \(a(S)\to a(G)\) est une équivalence faible. Donc \(a(S(G))\to a(G)\) l'est aussi.
\end{proofdel}

\section*{5 Deux foncteurs}

\subsection*{5.1 Le foncteur \(X\mapsto X/G\)}

Il y a un morphisme naturel de sites
\[
p:(QP/G)_{\Nis}\longrightarrow (QP)_{\Nis}
\]
donné par le foncteur
\[
p^*:\ QP\longrightarrow QP/G,\qquad X\longmapsto (X\text{ avec l'action triviale de }G).
\]
En effet \(p^*\) commute aux limites finies et transforme les familles couvrantes en familles couvrantes. En particulier il est continu: si \(F\) est un faisceau sur \((QP/G)_{\Nis}\), le préfaisceau
\[
X\longmapsto F(X\text{ avec action triviale de }G)
\]
est un faisceau sur \((QP)_{\Nis}\). Le foncteur \(p^*\) admet un adjoint à gauche
\[
q:\ QP/G\longrightarrow QP,\qquad X\longmapsto X/G .
\]
Comme \(p^*\) est continu, \(q\) est cocontinu, et l'image inverse \(p^*\) des faisceaux sur \((QP)_{\Nis}\) vers les faisceaux sur \((QP/G)_{\Nis}\) est
\[
F\longmapsto \bigl(\text{faisceau associé à }X\longmapsto F(X/G)\bigr).
\]

\begin{proposition*}[43]
Le foncteur cocontinu \(q:X\mapsto X/G\) est aussi continu. Autrement dit, si \(F\) est un faisceau sur \(QP\), le préfaisceau \(X\mapsto F(X/G)\) sur \(QP/G\) est un faisceau.
\end{proposition*}

\begin{proofdel}
Par le lemme 2, il suffit de tester la propriété de faisceau de \(X\mapsto F(X/G)\) pour un recouvrement de \(X\) obtenu par tiré en arrière d'un recouvrement de Nisnevich \(U_i\to X/G\). Le passage au quotient commute au changement de base plat. En prenant \(X/G\) pour base, cela identifie \(X\times_{X/G}U_i\to U_i\) à un quotient par \(G\) égal à \(U_i\). Il en va de même pour les intersections multiples. La propriété de faisceau pour le recouvrement de \(X\) par les \(X\times_{X/G}U_i\) se ramène donc à la propriété de faisceau de \(F\) pour le recouvrement \((U_i)\) de \(X/G\).
\end{proofdel}

Le foncteur \(q:X\mapsto X/G\) donne une paire de foncteurs adjoints entre préfaisceaux sur \(QP/G\) et \(QP\). Comme \(q\) est continu, il induit une paire analogue sur les faisceaux. Nous la notons
\[
(p_\#,p^*),
\]
de sorte que l'on dispose d'une suite d'adjonctions \((p_\#,p^*,p_*)\). Si \(F\) sur \(QP/G\) est représentable, \(F=h_X\), alors \(p_\#F=h_{X/G}\). En particulier \(p_\#\) transforme le faisceau final, ou point, sur \(QP/G\), en le faisceau final sur \(QP\), et \((p_\#,p^*)\) est aussi une paire adjointe dans la catégorie des faisceaux pointés. Il est clair que \(p_\#\) envoie les faisceaux solides sur des faisceaux solides. On a aussi le résultat suivant.

\begin{proposition*}[44]
Soit \(F\) un faisceau pointé solide relativement aux plongements ouverts \(U\subset X\) de schémas lisses tels que l'action de \(G\) sur \(X\) soit libre hors de \(U\). Alors \(p_\#(F)\) est solide relativement aux plongements ouverts de schémas lisses.
\end{proposition*}

\begin{proofdel}
Si \(X'\) est l'ouvert de \(X\) où l'action de \(G\) est libre, alors \(U\cup X'\to X\) appartient à la même classe, et, en posant \(U'=U\cap X'\), une somme amalgamée de \(U\to X\) est aussi une somme amalgamée de \(U'\to X'\). On s'est donc ramené au cas où l'action est libre partout. Puisque \(p_\#\) est adjoint à gauche, il préserve les colimites et en particulier les sommes amalgamées. Il envoie \(h_U\) sur \(h_{U/G}\). L'assertion se réduit donc au lemme suivant.
\end{proofdel}

\begin{lemma*}[17]
Si \(G\) agit librement sur \(U\), lisse sur \(S\), alors \(U/G\) est lisse.
\end{lemma*}

\begin{proofdel}
Si \(G\) est fini étale, par exemple \(S_n\), ce qui est le cas qui nous intéresse le plus, c'est clair puisque \(U\to U/G\) est étale. En général, l'hypothèse que \(G\) agit librement sur \(U\) implique que \(U\) est un \(G\)-torseur sur \(U/G\). En particulier \(U\to U/G\) est fidèlement plat. Comme \(U\) est plat sur \(S\), cela force \(U/G\) à être plat sur \(S\). Pour vérifier que \(U/G\) est lisse sur \(S\), il suffit de le vérifier fibre géométrique par fibre géométrique. Pour un point géométrique \(\bar s\) de \(S\), la lissité de \((U/G)_{\bar s}\) revient à la régularité. Comme \(U_{\bar s}\) est lisse sur \(\bar s\), donc régulier, et que \(U_{\bar s}\to(U/G)_{\bar s}\) est fidèlement plat, ceci résulte de [4], application du critère cohomologique de régularité de Serre.
\end{proofdel}

\begin{proposition*}[45]
Le foncteur \(p_\#\) respecte les équivalences locales, respectivement les \(\Aone\)-équivalences, entre faisceaux simpliciaux terme à terme ind-solides.
\end{proposition*}

\begin{proofdel}
Soit \(f:F\to F'\) une équivalence locale entre faisceaux simpliciaux terme à terme ind-solides sur \(QP/G\). Pour vérifier que \(p_\#(f)\) est une équivalence locale, il suffit de voir que pour tout \(X\in QP\) et tout \(x\in X\), le morphisme
\[
p_\#(F)(\Spec\mathcal O^h_{X,x})\longrightarrow
p_\#(F')(\Spec\mathcal O^h_{X,x})
\]
est une équivalence faible d'ensembles simpliciaux. Comme \(p_\#\) est un adjoint à gauche, le foncteur
\[
F\longmapsto p_\#(F)(\Spec\mathcal O^h_{X,x})
\tag{45.1}
\]
commute aux colimites. Pour un plongement ouvert \(U\to X\) dans \(QP/G\), le morphisme \(U/G\to X/G\) est encore un plongement ouvert, et le théorème 6 s'applique à (45.1).

Soit \(f:F\to F'\) une \(\Aone\)-équivalence. Considérons le carré
\[
\begin{tikzcd}[column sep=2.8em,row sep=1.8em]
S(F) \arrow[r] \arrow[d] & S(F') \arrow[d]\\
F \arrow[r] & F' .
\end{tikzcd}
\]
Par la première partie de la proposition, \(p_\#\) transforme les flèches verticales en équivalences locales. Comme \(p_\#\) commute aux colimites, le lemme 13 implique que \(p_\#(S(f))\) appartient à la \(\bar\Delta\)-clôture de la classe contenant les \(p_\#((K_Q)_+\to X_+)\), pour \(Q\) distingué supérieur, et les \(p_\#((X\times\Aone)_+\to X_+)\), pour \(X\in QP/G\). Par le théorème 4, il reste à prouver que les morphismes de ces deux types sont des \(\Aone\)-équivalences. Pour le premier type cela résulte de la première moitié de la proposition. Pour le second, cela résulte du fait que
\[
p_\#((X\times\Aone)_+)\longrightarrow p_\#(X_+)
\]
et
\[
p_\#(X_+)\times\Aone\longrightarrow p_\#((X\times\Aone)_+)
\]
sont des équivalences d'homotopie \(\Aone\) mutuellement inverses.
\end{proofdel}

On définit
\[
Lp_\#:\Ho_{A^1}(QP/G)\longrightarrow \Ho_{A^1}(QP)
\]
(et de même dans le cadre pointé) par
\[
Lp_\#(F):=p_\#(S(F)),
\]
où \(S(F)\) est défini dans la construction 29. La proposition 45 montre que \(Lp_\#\) est bien défini et que, pour \(F\) terme à terme ind-solide, on a \(Lp_\#(F)\simeq p_\#(F)\).

\subsection*{5.2 Le foncteur \(X\mapsto X^T\)}

Comme au paragraphe 3.1, nous fixons \(G\) et \(S\). Nous fixons aussi \(T\in QP/G\), fini et plat sur \(S\).

Pour un préfaisceau \(F\) sur \(QP/G\), on définit \(F^T\) comme l'objet Hom interne \(\underline{\Hom}(h_T,F)\). Sa valeur sur \(U\) est
\[
F^T(U)=F(U\times_S T).
\]
Si \(F\) est un faisceau, alors \(F^T\) en est un.

\begin{example*}[6]
Prenons \(G\) et \(T\) déduits du groupe fini \(S_n\) agissant sur \(\{1,\ldots,n\}\) par permutations. Si \(F\) est représenté par \(X\), avec action triviale de \(S_n\), alors \(F^T\) est représenté par \(X^n\), où \(S_n\) agit par permutation des facteurs.
\end{example*}

\begin{remark*}[1]
Si \(F\) est représentable, et représenté par \(X\) lisse sur \(S\), alors \(F^T\) est représentable, et représenté par un \(G\)-schéma lisse. Plus précisément, si \(F\) est représenté par \(X\in QP/G\), considérons le foncteur contravariant sur \(Sch/S\) des morphismes de schémas de \(T\) vers \(X\),
\[
U\longmapsto \Hom_S(T\times_S U,\; X\times_S U).
\]
Ce foncteur est représenté par un schéma \(Y\), quasi-projectif sur \(S\), et lisse si \(X\) l'est. Le schéma \(Y\) porte l'action évidente de \(G\), et \((Y,G)\) représente \(F^T\). En effet, en associant à un morphisme \(T\to X\) son graphe, on envoie le foncteur dans le foncteur de Hilbert des sous-schémas finis de \(T\times_S X\) de même rang que \(T\). Ce foncteur de Hilbert est représenté par un schéma de Hilbert quasi-projectif. La condition qu'un sous-schéma fini plat soit le graphe d'un morphisme de \(T\) vers \(X\) est ouverte; elle définit le sous-schéma ouvert \(Y\) requis. Si \(X\) est lisse, la lissité de \(Y\) résulte du critère infinitésimal de relèvement. L'action de \(G\) est évidente. Pour \(Z\in QP/G\), on a les identifications
\[
\Hom_{QP/G}(Z,Y)
\simeq \Hom_S(Z,\Hom(T,X))
\simeq \Hom_S(Z\times_S T,X)
\simeq F^T(Z).
\]
\end{remark*}

Soit \(\mathcal E\) une classe de plongements ouverts dans \((QP/G)_{\Nis}\). Nous dirons simplement «solide» pour \(\mathcal E\)-solide.

\begin{theorem*}[7]
Si \(F\) est un faisceau solide sur \((QP/G)_{\Nis}\), alors \(F^T\) est solide.
\end{theorem*}

Si \(A\to F\) est ouvert, c'est-à-dire relativement représentable par des plongements ouverts équivariants, il existe une suite naturelle de faisceaux intermédiaires entre \(A^T\) et \(F^T\). Dans la situation de l'exemple 6, et pour \(h_U\to h_X\), ils sont représentés par les sous-schémas ouverts équivariants de \(X^n\) formés des \(n\)-uplets \((x_1,\ldots,x_n)\) pour lesquels au moins \(i\) des \(x_j\) sont dans \(U\).

Formellement, une section de \(F^T\) sur \(Z\) est une section \(s\) de \(F\) sur \(T\times_S Z\). Soit \(U(s)\) l'ouvert équivariant de \(T\times_S Z\) sur lequel \(s\) appartient à \(A\). Le faisceau \((F,A)^T_i\) est le sous-faisceau de \(F^T\) formé des \(s\) tels que toutes les fibres de \(U(s)\to Z\) soient de degré au moins \(i\). La condition sur le degré étant ouverte, l'inclusion \((F,A)^T_i\subset F^T\) est ouverte. Pour \(i=0\), ce faisceau est \(F^T\); pour \(i\) assez grand, il est \(A^T\).

\begin{lemma*}[18]
Supposons que \(A\to F\) soit déduit par somme amalgamée d'un morphisme ouvert \(B\to C\), de sorte que
\[
\begin{tikzcd}[column sep=2.8em,row sep=1.8em]
B \arrow[r] \arrow[d] & C \arrow[d]\\
A \arrow[r] & F
\end{tikzcd}
\tag{45.2}
\]
soit cocartésien. Alors, pour tout \(i\), le carré cartésien
\[
\begin{tikzcd}[column sep=3.2em,row sep=1.8em]
\text{produit fibré} \arrow[r] \arrow[d] &
(A\amalg C,A)^T_i \arrow[d]\\
(F,A)^T_{i+1} \arrow[r] &
(F,A)^T_i
\end{tikzcd}
\tag{45.3}
\]
est aussi cocartésien.
\end{lemma*}

\begin{proofdel}
Le site \((QP/G)_{\Nis}\) possède assez de points. Par le lemme 2, pour chaque \(X\in QP/G\) et chaque \(x\in X/G\), le foncteur
\[
F\longmapsto \varinjlim F(X\times_{X/G}U),
\]
où \(U\) parcourt les voisinages de Nisnevich de \(x\) dans \(X/G\), est un point. Ces points sont représentés par des schémas henséliens \(G\)-locaux \(Z\): réunions disjointes finies de schémas locaux henséliens essentiellement de type fini sur \(S\), avec \(Z/G\) local.

Nous montrons que (45.3) devient cocartésien après application d'un tel foncteur fibre \(F\mapsto F(Z)\). Fixons \(s\in(F,A)^T_i(Z)\), vu comme une section de \(F\) sur \(T\times_S Z\), et soit \(U\subset T\times_S Z\) l'ouvert où \(s\) appartient à \(A\). L'hypothèse signifie que le degré de la fibre \(U_z\) est au moins \(i\), indépendamment du point fermé \(z\) de \(Z\). Il reste à montrer que, si \(s\) n'appartient pas à \((F,A)^T_{i+1}(Z)\), donc si ce degré vaut exactement \(i\), alors \(s\) admet un relèvement unique à \((A\amalg C,A)^T_i(Z)\).

Le schéma \(T\times_S Z\) est une réunion disjointe de schémas henséliens \(G\)-locaux. Sur chaque composante où \(s\) n'appartient pas à \(A\), la cocartésianité de (45.2) donne un relèvement unique à \(C\). Sur les composantes où \(s\) appartient à \(A\), le relèvement est la section donnée dans \(A\). On obtient ainsi une section de \(A\amalg C\) sur \(T\times_S Z\) relevant \(s\), et elle est unique par le même raisonnement composante par composante.
\end{proofdel}

\begin{proofnamed}{Preuve du théorème 7}
Comme \(F\) est solide, il apparaît au bout d'une suite
\[
\emptyset=F_0\to F_1\to\cdots\to F_n=F
\]
où chaque \(F_i\to F_{i+1}\) est une somme amalgamée d'un plongement ouvert dans \(QP/G\). On prouve par récurrence sur \(i\) que, pour tout \(Z\), \((F_i\amalg h_Z)^T\) est solide. Pour \(i=0\), \(F_0\) est représentable, et \((F_0\amalg h_Z)^T\) l'est aussi par l'énoncé de représentabilité ci-dessus. Pour l'étape de récurrence, on applique le lemme suivant à \(F_i\amalg h_Z\to F_{i+1}\amalg h_Z\).
\end{proofnamed}

\begin{lemma*}[19]
Soit
\[
\begin{tikzcd}[column sep=2.8em,row sep=1.8em]
h_U \arrow[r] \arrow[d] & h_X \arrow[d]\\
F \arrow[r] & G
\end{tikzcd}
\]
un carré cocartésien avec \(U\) ouvert dans \(X\). Supposons que, pour tout \(Z\), \((F\amalg h_Z)^T\) soit solide. Alors \(F^T\to G^T\) est solide.
\end{lemma*}

\begin{proofdel}
Comme \(F\to G\) est ouvert par la proposition 33, les \((G,F)^T_i\) sont définis. Il suffit de montrer que, pour chaque \(i\), le morphisme ouvert
\[
(G,F)^T_{i+1}\longrightarrow (G,F)^T_i
\]
est solide. Par le lemme 18, ce morphisme s'insère dans un carré à la fois cartésien et cocartésien
\[
\begin{tikzcd}[column sep=3.0em,row sep=1.8em]
\text{produit fibré} \arrow[r] \arrow[d] &
(F\amalg h_X,F)^T_i \arrow[d]\\
(G,F)^T_{i+1} \arrow[r] &
(G,F)^T_i .
\end{tikzcd}
\tag{45.4}
\]
Par hypothèse \((F\amalg h_X)^T\) est solide; donc \((F\amalg h_X,F)^T_i\) l'est aussi, par la proposition 37 appliquée au morphisme ouvert vers \((F\amalg h_X)^T\). La flèche verticale droite de (45.4) est ouverte, donc solide par la proposition 36, et la flèche horizontale inférieure est solide comme somme amalgamée d'un morphisme solide.
\end{proofdel}

\begin{example*}[7]
Il n'est pas toujours vrai que, si \(f:A\to F\) est un morphisme solide, alors \(f^T\) soit solide. Prenons \(G\) trivial et \(T\) formé de deux points, de sorte que \(F^T=F^2\). Pour tout faisceau \(F\), l'inclusion \(F\to F\amalg\pt\) est solide. En appliquant \((-)^T\), on obtient
\[
F^2\longrightarrow F^2\amalg F\amalg F\amalg \pt .
\]
En tirant en arrière par le morphisme naturel de \(F\) vers l'un des deux facteurs \(F\), morphisme ouvert, on voit que si \(f^T\) était solide, alors \(F\) serait solide.
\end{example*}

\begin{corollary*}[46]
Si \(f:F\to G\) est ouvert et si \(G\) est solide, alors
\[
f^T:F^T\longrightarrow G^T
\]
est solide. En particulier, si \(G\) est solide pointé, alors \(G^T\) est solide pointé.
\end{corollary*}

\begin{proofdel}
Le morphisme \(f^T\) est ouvert, et l'on applique le théorème 7 et la proposition 36.
\end{proofdel}

Nous définissons maintenant, pour les faisceaux pointés sur \((QP/G)_{\Nis}\), une variante smash de la construction \(F\mapsto F^T\). Si le point marqué \(\pt\to F\) est ouvert, c'est
\[
F^{\wedge T}:=F^T/(F,\pt)^T_1,
\]
c'est-à-dire \(F^T\) avec \((F,\pt)^T_1\) contracté au nouveau point-base. Cette définition se répète pour tout faisceau pointé si, pour tout monomorphisme \(A\to F\), on définit \((F,A)^T_1\) comme le sous-faisceau suivant de \(F^T\): une section \(s\) de \(F^T(X)=F(X\times T)\) est dans \((F,A)^T_1\) si, pour tout \(X'\to X\) non vide, il existe un diagramme commutatif
\[
\begin{tikzcd}[column sep=2.8em,row sep=1.7em]
X'' \arrow[r] \arrow[d] & X\times T \arrow[d]\\
X' \arrow[r] & X
\end{tikzcd}
\]
avec \(X''\) non vide et \(s\) appartenant à \(A\) sur \(X''\).

\begin{example*}[8]
Sous les hypothèses de l'exemple 6, si \(U\) est ouvert dans \(X\), de complément \(Z\), et si \(F=h_X/h_U\), alors \(F^{\wedge T}\) est \(h_{X^n}/h_{X^n-Z^n}\), où \(X^n\) porte son action naturelle de \(S_n\). En particulier, si \(F\) est l'espace de Thom d'un fibré vectoriel \(E\) sur \(Y\), alors \(F^{\wedge T}\) est l'espace de Thom du fibré vectoriel \(S_n\)-équivariant \(\bigoplus_1^n E\) sur \(Y^n\).
\end{example*}

\begin{proposition*}[47]
Si un faisceau pointé \(F\) est solide pointé, c'est-à-dire si \(\pt\to F\) est solide, alors \(F^{\wedge T}\) est solide pointé.
\end{proposition*}

\begin{proofdel}
Comme \(F^T\) est solide, le morphisme ouvert \((F,\pt)^T_1\to F^T\) est solide par la proposition 36, et \(\pt\to F^{\wedge T}\) s'en déduit par somme amalgamée.
\end{proofdel}

La définition de \(F^{\wedge T}\) implique immédiatement le résultat suivant.

\begin{lemma*}[20]
Soit \(Z\) un schéma hensélien \(G\)-local. Alors
\[
F^{\wedge T}(Z)=\bigwedge_i F((T\times Z)_i),
\]
où les \((T\times Z)_i\) sont les composantes henséliennes \(G\)-locales de \(T\times Z\).
\end{lemma*}

\begin{proposition*}[48]
Le foncteur \(F\mapsto F^{\wedge T}\) respecte les équivalences locales et les \(\Aone\)-équivalences.
\end{proposition*}

\begin{proofdel}
Soit \(f:F\to G\) une équivalence locale. Pour vérifier que \(F^{\wedge T}\to G^{\wedge T}\) est une équivalence locale, il suffit de tester sur tout schéma hensélien \(G\)-local \(Z\). Par le lemme 20, le morphisme devient le smash-produit des morphismes
\[
F((T\times Z)_i)\longrightarrow G((T\times Z)_i),
\]
donc une équivalence faible d'ensembles simpliciaux.

Supposons maintenant que \(f\) soit une \(\Aone\)-équivalence. Par la première partie il suffit de montrer que
\[
S(f)^{\wedge T}:S(F)^{\wedge T}\longrightarrow S(G)^{\wedge T}
\]
est une \(\Aone\)-équivalence. On utilise la caractérisation des \(\Aone\)-équivalences donnée au théorème 5. Comme \(F\mapsto F^{\wedge T}\) commute aux colimites filtrantes et préserve les équivalences locales, il suffit de vérifier qu'il transforme les équivalences d'homotopie \(\Aone\) en équivalences d'homotopie \(\Aone\). Cela résulte du morphisme naturel
\[
F^{\wedge T}\wedge(\Aone)_+\longrightarrow (F\wedge(\Aone)_+)^{\wedge T}.
\]
\end{proofdel}

\section*{Remerciements}

Travail soutenu par les subventions NSF DMS-97-29992 et DMS-9901219, ainsi que par la Fondation Ambrose Monell.

\section*{Références}

\begin{enumerate}[label={[\arabic*]}]
\item K. S. Brown et S. M. Gersten, \emph{Algebraic \(K\)-theory and generalized sheaf cohomology}, Lecture Notes in Mathematics 341, Springer, 1973, p. 266--292.
\item A. Grothendieck, M. Artin et J.-L. Verdier, \emph{Théorie des topos et cohomologie étale des schémas} (SGA 4), Lecture Notes in Mathematics 269, 270, 305, Springer, 1972--1973.
\item A. Grothendieck et M. Demazure, \emph{Schémas en groupes} (SGA 3), Lecture Notes in Mathematics 151, 152, 153, Springer, 1970.
\item A. Grothendieck et J. Dieudonné, \emph{Étude locale des schémas et des morphismes de schémas} (EGA 4), Publ. Math. IHES 20, 24, 28, 32, 1964--1967.
\item Mark Hovey, \emph{Model Categories}, American Mathematical Society, Providence, 1999.
\item S. Mac Lane, \emph{Categories for the Working Mathematician}, Graduate Texts in Mathematics 5, Springer, 1971.
\item Carlo Mazza, Vladimir Voevodsky et Charles Weibel, \emph{Lecture Notes on Motivic Cohomology}, Clay Mathematics Monographs 2, AMS, Providence, RI, 2006.
\item Fabien Morel et Vladimir Voevodsky, \emph{\(\Aone\)-homotopy theory of schemes}, Publ. Math. IHES 90 (1999), 45--143.
\item Charles A. Weibel, \emph{An Introduction to Homological Algebra}, Cambridge University Press, 1994.
\item P. Gabriel et M. Zisman, \emph{Calculus of Fractions and Homotopy Theory}, Springer-Verlag, Berlin, 1967.
\item J. F. Jardine, \emph{Simplicial presheaves}, J. Pure Appl. Algebra 47 (1987), 35--87.
\end{enumerate}
\end{document}
